\chapter{What a Verified Notebook Is}\label{ch:what}

\begin{quote}\itshape
``The answer is $2x\cos x + x^2\sin x$.''\;---\;``Says who?''
\end{quote}

A computer-algebra system prints an answer and asks to be believed. Usually it is right, and
usually that is enough. This book is about a notebook built on a different bargain: it prints the
answer \emph{and} the steps, and for each step it says what would have to be true for the step
to be a lie. Sometimes nothing: the step is a theorem. Sometimes one named hypothesis. Sometimes
``this step is a guess, and the step below it checks the guess''. And sometimes ``no theorem yet''.
The notebook never says more than it can back.

\section{Show your work}

Type \texttt{diff(x\^{}2 * sin(x), x)} into the notebook and it answers
$2x\sin x + x^2\cos x$, with seven steps underneath: the product rule, the power rule, a constant
folded, an identity applied, the chain rule, an identity applied again. Each step is a line, the
whole term after the step, with the part that changed underlined. Click any subterm of any line
--- the $\cos x$ in the answer, say --- and a panel names the step that \emph{created} it (the
chain rule), the steps that merely \emph{copied} it along, and the steps that fired \emph{inside}
it, together with the proof status of each. That panel is Chapter~\ref{ch:origins}.

The four statuses are the book's spine, so they are worth stating once, properly.

\begin{defbox}[The four statuses]
  Every rule the engine can fire is reported to the notebook with one of four labels.
  \begin{description}[nosep]
    \item[\status{verified}] An unconditional theorem states that the output of the rule means the
      same as its input, in the semantics of the chapter that proved it (the integers, the reals,
      the reals under a derivative, the rationals in a linear system, \dots).
    \item[\status{conditional}] The theorem needs a side condition --- \emph{and the side condition
      is shown to be necessary}: a counterexample is proved, not just suspected. The notebook shows
      the condition.
    \item[\status{checked}] The step is a guess. A later step verifies the guess by an independent
      computation whose own steps carry their own statuses. (Integration, Chapter~\ref{ch:integrate}.)
    \item[\status{unverified}] There is no theorem yet. The notebook says so, in those words.
  \end{description}
  A step with nested steps inherits the weakest status below it.
\end{defbox}

The distinction between the first two is the one that gives the book most of its stories. A
computer-algebra system simplifies $x/x$ to $1$. So does this engine. But $x/x$ is not $1$ at
$x = 0$, and rather than pretend otherwise, or refuse the simplification everyone expects, the
engine keeps the rule and \emph{writes down the assumption}: the rule is \status{conditional} on
the merged base being nonzero, and a theorem in Chapter~\ref{ch:sound} proves that without the
assumption the rule turns $0$ into $1$. That is the whole attitude of the project in one rule.

\section{One engine}

There is exactly one implementation of the mathematics, and it is written in Lean~4. It compiles
to a native executable (behind an HTTP host, or as a command-line tool speaking JSON-RPC on
standard input) and, from the same source, to WebAssembly running in a web worker. The notebook
is a static page; it draws terms with KaTeX and draws plots and Hasse diagrams as SVG, and it never
evaluates anything. Every number, every rewrite, every ``this is a lattice'' comes over the wire
from the engine, with the derivation attached.

\begin{center}
\resizebox{\linewidth}{!}{\begin{tikzpicture}[node distance=9mm]
  \node[box] (nb) {notebook (static page)\\KaTeX, SVG, no mathematics};
  \node[box, right=9mm of nb] (proto) {JSON-RPC over a Transport\\worker · HTTP · stdio};
  \node[box, right=9mm of proto] (eng) {the engine, Lean 4\\native and wasm};
  \node[box, below=of eng] (proofs) {\file{proofs/}: theorems about\\the engine, with Mathlib};
  \draw[edge] (nb) -- (proto);
  \draw[edge] (proto) -- (eng);
  \draw[edge, dashed] (proofs) -- node[right, font=\scriptsize] {imports} (eng);
\end{tikzpicture}}
\end{center}

Two decisions from the first week held for the whole project.

\paragraph{Two packages.} The engine, in \file{engine/}, imports nothing but Lean's prelude
\lc{Init}. The reason is mechanical: the WebAssembly build emits C for every imported Lean module,
and \lc{Init} alone is 631 of them; Mathlib would be tens of thousands. So the real numbers,
which need Mathlib, cannot live in the engine. They live in a second package, \file{proofs/},
which imports the engine and Mathlib and contains theorems only. A script,
\file{scripts/check-engine-deps.sh}, fails the build if the engine ever acquires a dependency.
The consequence for the reader: every theorem in this book either lives in the engine (and is
proved with Lean's core tactics --- \lc{omega}, \lc{simp}, \lc{grind}, \lc{decide}) or lives in
\file{proofs/} and may use anything Mathlib knows about $\mathbb{R}$. The book says which.

\paragraph{The protocol is the seam.} The notebook and the engine share a small JSON-RPC contract
(\file{packages/protocol}). Its one design rule --- rule~5 --- is that new capabilities are added
as \emph{optional fields and optional methods}, never as changes. Every chapter of this book that
added something to the wire (paths, proof statuses, origin traces, the de~Bruijn view, the Hasse
diagram) added it as a field a frontend may ignore. The notebook from week one would still run
against the engine from the last chapter.

\section{The spike: does the engine reach the browser?}

Before any mathematics, one question decided whether the project was possible: can a Lean~4
program of this size run in a browser, on \emph{current} Lean, fast enough to answer as you type?

\begin{logbox}[Milestone 0]
  Lean stopped shipping a prebuilt 32-bit WebAssembly runtime after v4.15. The project's policy is
  the latest stable toolchain (v4.33.1 at the time of writing), so the runtime and \lc{Init} were
  built from source for \texttt{wasm32} with Emscripten: about nine minutes the first time, cached
  afterwards. The result was a 1.4~MB \texttt{engine-lean.wasm} that answered its first request
  96~ms after the worker started, and evaluated \texttt{x + 0} in 50~µs per call after warm-up. No
  threads, so no cross-origin isolation headers: the notebook stays a plain static directory.

  One upstream bug surfaced on the way. On a 32-bit target, the runtime's conversion of a string
  to a list of characters built the empty list with a boxed \lc{UInt32} zero, which on
  \texttt{wasm32} is a heap object, not a scalar; the patch is three lines and lives in
  \file{engine/wasm/}. It looked like a reference-counting error for an hour.
\end{logbox}

The spike settled the shape of everything after it. The engine would be one program; the
notebook would be thin; the proofs would live next door; and ``verified'' would mean what the
statuses say. The rest of the book is what it took to make the statuses true.

\begin{trybox}[Your first three cells]
\begin{minted}{text}
diff(x^2 * sin(x), x)
rref([1,2,3;4,5,6;7,8,10])
sqrt(50) - sqrt(18)
\end{minted}
  Open the work on each, and click a subterm of an answer. The third cell's answer is $2\sqrt{2}$,
  and Chapter~\ref{ch:cannot} explains why the engine could not write it that way as a term.
\end{trybox}
